\documentclass[../main.tex]{subfiles}
\begin{document}

\section{Kết luận}
\label{sec:ket-luan}

Đồ án đã hoàn thành mục tiêu xây dựng một hệ thống điểm danh thông minh, chống gian lận trên
nền tảng Zalo Mini App, áp dụng được cho môi trường đào tạo đại học. Đối chiếu với các mục
tiêu cụ thể đặt ra ở mục~\ref{sec:muc-tieu}, kết quả đạt được như sau:

\begin{itemize}[leftmargin=1.5em]
  \item \textbf{Mã QR động chống chia sẻ.} Hệ thống đã hiện thực mã QR ký bằng HMAC-SHA256,
  xoay theo chu kỳ 30 giây và đồng bộ giữa điện thoại giảng viên với màn hình máy chiếu, qua
  đó vô hiệu hoá việc chụp và chia sẻ lại mã.

  \item \textbf{Xác minh khuôn mặt phía thiết bị.} Việc nhận diện khuôn mặt được thực hiện
  ngay trên thiết bị người dùng, chỉ lưu véc-tơ đặc trưng 128 chiều thay vì ảnh gốc.

  \item \textbf{Xác minh ngang hàng.} Cơ chế mỗi sinh viên tự quét mã của bạn (tối thiểu một
  lượt, chỉ lượt tự quét mới được tính) đã tạo ra bằng chứng về sự hiện diện vật lý, chống điểm danh từ xa.

  \item \textbf{Trang quản trị web.} Đã xây dựng trang quản trị cho giảng viên và phòng đào
  tạo với các chức năng quản lý lớp, giám sát phiên theo thời gian thực và xuất báo cáo Excel.

  \item \textbf{Phát hiện gian lận nhiều lớp.} Hệ thống tổng hợp bốn yếu tố xác minh thành
  điểm tin cậy và bổ sung hàm phân tích phía máy chủ để phát hiện, lập báo cáo tự động các
  mẫu hành vi đáng ngờ.
\end{itemize}

Về sản phẩm, đồ án đã tạo ra ba thành phần hoạt động được: một Mini App cho sinh viên và giảng
viên, một trang quản trị web, và một tập các hàm xử lý phía máy chủ trên nền tảng Firebase
(chi tiết ở Chương~\ref{chapter:Experiment}). Hệ thống đã được kiểm thử theo bộ 72 kịch bản và
đạt mức điểm tự đánh giá bảo mật 7{,}6/10. Những đóng góp nổi bật --- đặc biệt là cơ chế xác
minh ngang hàng và mã QR động có đồng bộ đa thiết bị
(Chương~\ref{chapter:SolutionAndContribution}) --- cho thấy có thể nâng cao đáng kể tính trung
thực của điểm danh bằng các lớp bảo vệ chi phí thấp, chỉ dựa trên thiết bị sẵn có của người
dùng mà không cần phần cứng chuyên dụng.

\section{Hạn chế}
\label{sec:han-che}

Bên cạnh các kết quả đạt được, đồ án vẫn còn một số hạn chế cần được nhìn nhận một cách thẳng
thắn:

\begin{itemize}[leftmargin=1.5em]
  \item \textbf{Phạm vi triển khai.} Hệ thống mới được thử nghiệm ở mức thử nghiệm (pilot)
  trong khuôn khổ một đồ án, chưa triển khai ở quy mô toàn trường; kết quả kiểm thử do đó còn
  mang tính định tính, chưa có số liệu trên tập người dùng lớn.

  \item \textbf{Một số nội dung ngoài phạm vi.} Hệ thống chưa tích hợp với hệ thống quản lý
  học tập (LMS) của nhà trường, chưa hỗ trợ thẻ sinh viên RFID/NFC, và mới dừng ở mức phát
  hiện sự sống (liveness) cơ bản nên vẫn còn rủi ro bị qua mặt bằng ảnh hoặc video.

  \item \textbf{Bảo mật và vận hành.} Một số dữ liệu phía thiết bị (\texttt{localStorage})
  chưa được mã hoá; cơ chế giới hạn tần suất và chống phát lại còn lưu trong bộ nhớ tạm nên
  không bền vững khi mở rộng quy mô; hệ thống cũng chưa có nhật ký kiểm toán và hiện vẫn vận
  hành trên gói dịch vụ miễn phí (xem mục~\ref{sec:bao-mat}).

  \item \textbf{Gian lận còn lại.} Xác minh ngang hàng vẫn có thể bị một nhóm sinh viên thông
  đồng cùng có mặt để xác minh khống cho nhau; lớp phân tích hành vi chỉ giảm thiểu chứ chưa
  loại bỏ hoàn toàn rủi ro này.

  \item \textbf{Hiệu năng.} Đồ án chưa thực hiện đo đạc hiệu năng định lượng (thời gian phản
  hồi, thông lượng) trên nhiều cấu hình thiết bị khác nhau.
\end{itemize}

\section{Hướng phát triển}
\label{sec:huong-phat-trien}

Từ những hạn chế nêu trên, đồ án đề xuất các hướng phát triển tiếp theo, sắp xếp theo mức độ
ưu tiên:

\begin{itemize}[leftmargin=1.5em]
  \item \textbf{Củng cố bảo mật và vận hành.} Mã hoá dữ liệu trên thiết bị; chuyển cơ chế giới
  hạn tần suất và chống phát lại sang lưu trữ bền vững; bổ sung nhật ký kiểm toán, tự đăng xuất
  sau thời gian dài và nâng cấp lên gói dịch vụ trả phí để khai thác đầy đủ năng lực máy chủ.

  \item \textbf{Nâng cao chống gian lận.} Áp dụng học máy để phát hiện gian lận dựa trên dữ
  liệu thay vì chỉ dựa trên luật; tăng cường phát hiện sự sống (chống giả mạo bằng ảnh, video)
  và nhận diện thiết bị bất thường.

  \item \textbf{Tích hợp hệ thống nhà trường.} Kết nối với hệ thống quản lý học tập (LMS), hỗ
  trợ thẻ sinh viên RFID/NFC và đăng nhập một lần (SSO) với tài khoản của trường.

  \item \textbf{Bảo vệ dữ liệu cá nhân.} Hướng tới mã hoá đầu cuối cho dữ liệu sinh trắc học,
  giảm thiểu lưu trữ thông tin định danh không cần thiết, và mời kiểm thử xâm nhập (penetration
  testing) độc lập.

  \item \textbf{Mở rộng sản phẩm.} Nhờ đã được điều chỉnh theo hướng trung lập về thương hiệu,
  hệ thống có thể được triển khai cho nhiều cơ sở đào tạo khác nhau, kèm theo các tính năng
  thống kê và thông báo phong phú hơn.
\end{itemize}

Tóm lại, đồ án đã xây dựng thành công một hệ thống điểm danh chống gian lận hoàn chỉnh và khả
thi về mặt triển khai, đồng thời chỉ ra một lộ trình rõ ràng để tiếp tục hoàn thiện hệ thống
thành một sản phẩm sẵn sàng cho môi trường thực tế.

\end{document}
